\documentclass[11pt]{article}
\usepackage[margin=1in]{geometry}
\usepackage{enumitem}
\usepackage{fancyhdr}
\usepackage{titlesec}
\usepackage{amsmath}
\usepackage{amssymb}

\pagestyle{fancy}
\fancyhf{}
\rhead{GARP RAI Practice Exam — Modules 3, 4, 5}
\lhead{100 Balanced Exam-Style Questions}
\cfoot{\thepage}

\titleformat{\section}{\large\bfseries}{\thesection.}{0.5em}{}
\titleformat{\subsection}{\normalsize\bfseries}{}{0em}{}

\title{\textbf{Risk and AI (RAI) Practice Exam}\\
	\large 100 Questions — Modules 3, 4, and 5\\
	\normalsize Balanced Options with Full Explanations}
\author{Original Practice Set}
\date{\today}

\begin{document}
	\maketitle
	
	\section*{Instructions}
	Each question has four choices (A--D). Select the best answer.
	Options are balanced in length and grammatical structure so that length
	is not a cue to the correct answer. Suggested time: 2.5 minutes per question.
	
	\newpage
	\section{Module 3: Fairness, Explainability, and Human Risk}
	
	\subsection*{Question 1}
	A medical AI diagnoses a rare disease affecting 0.1\% of the population.
	Untreated, the disease is almost always fatal. Which metric best assesses
	the algorithm's performance?
	
	A. Sensitivity, which measures the proportion of actual positives correctly identified.\\
	B. Precision, which measures the proportion of predicted positives that are correct.\\
	C. Accuracy, which measures the proportion of all predictions that are correct.\\
	D. Specificity, which measures the proportion of actual negatives correctly identified.
	
	\vspace{0.4em}
	\textbf{Correct Answer: A}
	
	\textbf{Explanation:}
	A is correct. Sensitivity (true positive rate) captures the proportion of
	actual positive cases correctly identified. Missing a true case is fatal.
	B is incorrect. Precision ignores false negatives, which are the critical
	error in this scenario.
	C is incorrect. Accuracy is misleading with rare events because a
	never-positive classifier scores 99.9\%.
	D is incorrect. Specificity captures true negatives, which matter less
	when the disease is curable with few side effects.
	
	\textbf{Learning Objective:} Describe various measures of group fairness.
	\textbf{Reference:} Module 3
	
	% =====================================================
	\subsection*{Question 2}
	How does demographic parity contribute to group fairness?
	
	A. It ensures that model accuracy is equal across all subgroups.\\
	B. It ensures that predicted positive rates are equal across subgroups.\\
	C. It ensures that model sensitivity is equal across all subgroups.\\
	D. It ensures that model precision is equal across all subgroups.
	
	\vspace{0.4em}
	\textbf{Correct Answer: B}
	
	\textbf{Explanation:}
	A is incorrect. Demographic parity says nothing about accuracy.
	B is correct. Demographic parity requires equal predicted positive rates.
	C is incorrect. Equal sensitivity is a different metric (equal opportunity).
	D is incorrect. Precision equality is not required by demographic parity.
	
	\textbf{Learning Objective:} Describe various measures of group fairness.
	\textbf{Reference:} Module 3
	
	% =====================================================
	\subsection*{Question 3}
	A fraud model is trained on 10{,}000 transactions, only 20 of which are
	fraudulent. Which form of algorithmic bias is most likely to occur?
	
	A. Measurement bias, arising from inconsistent feature measurement.\\
	B. Representation bias, arising from poor subgroup representation.\\
	C. Data composition bias, arising from severe class imbalance.\\
	D. Historical bias, arising from outdated training patterns.
	
	\vspace{0.4em}
	\textbf{Correct Answer: C}
	
	\textbf{Explanation:}
	A is incorrect. Measurement bias concerns how features are measured.
	B is incorrect. Representation bias concerns subgroup under-representation.
	C is correct. Severe class imbalance (20 vs.\ 9{,}980) is composition bias.
	D is incorrect. Historical bias concerns outdated data, not imbalance.
	
	\textbf{Learning Objective:} Describe sources of algorithmic bias.
	\textbf{Reference:} Module 3
	
	% =====================================================
	\subsection*{Question 4}
	A hiring model uses education level. To prevent education from
	influencing predictions, the best approach is to:
	
	A. Remove the education variable entirely from the trained model.\\
	B. Replace education with its average value during the scoring step.\\
	C. Add interaction terms between education and other model features.\\
	D. Replace education with binary dummy variables for each degree.
	
	\vspace{0.4em}
	\textbf{Correct Answer: B}
	
	\textbf{Explanation:}
	A is incorrect. Removing the variable leaves its influence embedded in
	correlated features.
	B is correct. Setting education to its average equalizes treatment at
	prediction time.
	C and D are incorrect. They change representation but do not remove
	influence.
	
	\textbf{Learning Objective:} Describe sources of algorithmic bias.
	\textbf{Reference:} Module 3, Section 2.3.1
	
	% =====================================================
	\subsection*{Question 5}
	A bank's chatbot sometimes gives incorrect information. The most
	appropriate action is:
	
	A. Collect additional customer inputs before generating each answer.\\
	B. Develop an ethical framework and periodically review outputs.\\
	C. Generate periodic risk reports summarizing chatbot performance.\\
	D. Implement design mechanisms requiring human verification of outputs.
	
	\vspace{0.4em}
	\textbf{Correct Answer: D}
	
	\textbf{Explanation:}
	A is incorrect. More inputs do not prevent hallucination.
	B is incorrect. A framework alone does not fix incorrect responses.
	C is incorrect. Reporting describes but does not mitigate the problem.
	D is correct. Human verification prevents incorrect outputs from reaching
	customers.
	
	\textbf{Learning Objective:} Discuss risks posed by AI to human autonomy.
	\textbf{Reference:} Module 3
	
	% =====================================================
	\subsection*{Question 6}
	To increase a fraud model's explainability, a credit card company should:
	
	A. Provide clear reasoning for the decisions produced by the model.\\
	B. Use a deep neural network to improve predictive performance.\\
	C. Provide access to the source code, data, and design rationale.\\
	D. Design the model with simple features and linear relationships.
	
	\vspace{0.4em}
	\textbf{Correct Answer: A}
	
	\textbf{Explanation:}
	A is correct. Explainability = understandable reasoning for decisions.
	B is incorrect. Deep networks reduce explainability.
	C is incorrect. Code/data access increases transparency, not
	explainability.
	D is incorrect. Simplicity increases interpretability, not
	explainability.
	
	\textbf{Learning Objective:} Describe explainability, interpretability,
	and transparency.
	\textbf{Reference:} Module 3
	
	% =====================================================
	\subsection*{Question 7}
	Customers complain they do not understand why loan applications are
	rejected. The source of reputational risk and its mitigation is:
	
	A. A privacy breach, mitigated by securing customer data access.\\
	B. Algorithmic bias, mitigated by using more diverse training data.\\
	C. Lack of explainability, mitigated by providing rejection reasons.\\
	D. Lack of transparency, mitigated by disclosing all training data.
	
	\vspace{0.4em}
	\textbf{Correct Answer: C}
	
	\textbf{Explanation:}
	A and B are incorrect. Neither is described.
	C is correct. The complaint is about unexplained rejections.
	D is incorrect. Disclosing training data is inappropriate and
	insufficient.
	
	\textbf{Learning Objective:} Describe sources of AI-related reputational
	risk.
	\textbf{Reference:} Module 3
	
	% =====================================================
	\subsection*{Question 8}
	A technology company wants to minimize data-related risk. The best step
	is:
	
	A. Retain all customer data regardless of sensitivity level.\\
	B. Give customers the ability to transfer or delete their data.\\
	C. Store customer data with multiple third-party vendors.\\
	D. Maintain an unencrypted original copy in a secure database.
	
	\vspace{0.4em}
	\textbf{Correct Answer: B}
	
	\textbf{Explanation:}
	A, C, D are incorrect. These increase data risk.
	B is correct. Data subject rights (transfer, deletion) reduce abuse.
	
	\textbf{Learning Objective:} Describe ethical principles of privacy.
	\textbf{Reference:} Module 4
	
	% =====================================================
	\subsection*{Question 9}
	Which best distinguishes explainability from interpretability?
	
	A. Explainability and interpretability are essentially identical.\\
	B. Interpretability concerns model structure; explainability concerns
	individual predictions.\\
	C. Explainability applies only to linear models; interpretability is
	universal.\\
	D. Interpretability requires deep networks; explainability requires
	shallow models.
	
	\vspace{0.4em}
	\textbf{Correct Answer: B}
	
	\textbf{Explanation:}
	A is incorrect. They are related but distinct.
	B is correct. Interpretability = model design; explainability =
	prediction explanation.
	C and D are incorrect. They reverse the actual relationship.
	
	\textbf{Learning Objective:} Describe explainability, interpretability,
	and transparency.
	\textbf{Reference:} Module 3
	
	% =====================================================
	\subsection*{Question 10}
	A model trained on majority-group data performs worse for a minority
	group. This is primarily:
	
	A. Measurement bias, caused by inconsistent feature collection.\\
	B. Representation bias, caused by subgroup under-representation.\\
	C. Historical bias, caused by outdated training patterns.\\
	D. Label bias, caused by incorrect target labeling.
	
	\vspace{0.4em}
	\textbf{Correct Answer: B}
	
	\textbf{Explanation:}
	A is incorrect. Measurement bias concerns feature measurement.
	B is correct. Under-representation causes poor subgroup performance.
	C and D are incorrect. Not indicated by the scenario.
	
	\textbf{Learning Objective:} Discuss sources of algorithmic bias.
	\textbf{Reference:} Module 4
	
	% =====================================================
	\subsection*{Question 11}
	Which action best supports procedural fairness in an AI lending system?
	
	A. Publishing the model's source code for public inspection.\\
	B. Providing a mechanism for applicants to contest decisions.\\
	C. Increasing the model's complexity to improve accuracy.\\
	D. Removing all applicant data after a decision is rendered.
	
	\vspace{0.4em}
	\textbf{Correct Answer: B}
	
	\textbf{Explanation:}
	A is incorrect. Code publication supports transparency.
	B is correct. Contest mechanisms support procedural fairness.
	C and D are incorrect. Unrelated to procedural fairness.
	
	\textbf{Learning Objective:} Discuss the principles of justice.
	\textbf{Reference:} Module 4
	
	% =====================================================
	\subsection*{Question 12}
	A model's decision threshold is lowered to reduce false negatives. This
	will most likely:
	
	A. Increase precision by filtering out more false positives.\\
	B. Increase recall by capturing more actual positive cases.\\
	C. Decrease sensitivity by reducing the true positive rate.\\
	D. Decrease recall by missing more actual positive cases.
	
	\vspace{0.4em}
	\textbf{Correct Answer: B}
	
	\textbf{Explanation:}
	A is incorrect. Lowering the threshold usually reduces precision.
	B is correct. More positives captured = higher recall.
	C and D are incorrect. Recall and sensitivity both increase.
	
	\textbf{Learning Objective:} Evaluate classification model performance.
	\textbf{Reference:} Module 2, Chapter 8
	
	% =====================================================
	\subsection*{Question 13}
	A bank's AI rejects a loan without giving specific reasons. Which
	principle is most directly violated?
	
	A. Autonomy, which concerns respecting individual decision-making.\\
	B. Non-maleficence, which concerns avoiding harm to others.\\
	C. Explainability, which concerns providing reasons for decisions.\\
	D. Beneficence, which concerns acting in the interest of others.
	
	\vspace{0.4em}
	\textbf{Correct Answer: C}
	
	\textbf{Explanation:}
	A, B, D are incorrect. Secondary principles in this scenario.
	C is correct. Failure to explain a decision is an explainability failure.
	
	\textbf{Learning Objective:} Discuss the principles of AI ethics.
	\textbf{Reference:} Module 4
	
	% =====================================================
	\subsection*{Question 14}
	``Model interpretability'' is most important when:
	
	A. The model is used internally for exploratory data analysis.\\
	B. Decisions affect individuals and require formal justification.\\
	C. The model is trained primarily on synthetic or generated data.\\
	D. The model has no exposure to regulatory or legal review.
	
	\vspace{0.4em}
	\textbf{Correct Answer: B}
	
	\textbf{Explanation:}
	A, C, D are incorrect. These contexts lower the need for
	interpretability.
	B is correct. High-stakes individual decisions require justification.
	
	\textbf{Learning Objective:} Describe explainability, interpretability,
	and transparency.
	\textbf{Reference:} Module 3
	
	% =====================================================
	\subsection*{Question 15}
	Which is the best example of a human-in-the-loop safeguard?
	
	A. Fully automated loan approval with periodic post-hoc audits.\\
	B. Human review of AI-flagged AML alerts before regulatory filing.\\
	C. Fully automated trading with post-trade reconciliation checks.\\
	D. Auto-generation of regulatory filings with quarterly review.
	
	\vspace{0.4em}
	\textbf{Correct Answer: B}
	
	\textbf{Explanation:}
	A, C, D are incorrect. These are largely automated processes.
	B is correct. Human review before consequential action = HITL.
	
	\textbf{Learning Objective:} Discuss risks posed by AI to human autonomy.
	\textbf{Reference:} Module 3
	
	% =====================================================
	\subsection*{Question 16}
	A face recognition system performs poorly on darker skin tones. This is
	most likely caused by:
	
	A. Sampling bias, resulting from under-representation in training data.\\
	B. Excessive model capacity, resulting in overfitting to fine details.\\
	C. Over-regularization, resulting in loss of model flexibility.\\
	D. Insufficient feature engineering, resulting in poor feature maps.
	
	\vspace{0.4em}
	\textbf{Correct Answer: A}
	
	\textbf{Explanation:}
	A is correct. Subgroup under-representation causes subgroup failure.
	B, C, D are incorrect. Not the primary cause.
	
	\textbf{Learning Objective:} Discuss sources of algorithmic bias.
	\textbf{Reference:} Module 4
	
	% =====================================================
	\subsection*{Question 17}
	Which is the most appropriate response when an AI hiring model shows
	disparate impact?
	
	A. Ignore the finding because overall model accuracy is high.\\
	B. Audit features, retrain with balanced data, and document findings.\\
	C. Deploy the model and monitor subgroup performance annually.\\
	D. Replace the model with a purely random candidate selection.
	
	\vspace{0.4em}
	\textbf{Correct Answer: B}
	
	\textbf{Explanation:}
	A is incorrect. Overall accuracy does not excuse disparate impact.
	B is correct. Audit, retrain, and document.
	C is incorrect. Annual monitoring is too slow for an active issue.
	D is incorrect. Random selection destroys utility.
	
	\textbf{Learning Objective:} Discuss sources of algorithmic bias.
	\textbf{Reference:} Module 4
	
	% =====================================================
	\subsection*{Question 18}
	A loan model uses zip code as a feature, producing discriminatory
	outcomes. This illustrates:
	
	A. Proxy discrimination, where a neutral feature encodes a protected one.\\
	B. Concept drift, where the input-output relationship changes over time.\\
	C. Overfitting, where the model memorizes noise in the training data.\\
	D. Data drift, where the distribution of input features changes over time.
	
	\vspace{0.4em}
	\textbf{Correct Answer: A}
	
	\textbf{Explanation:}
	A is correct. Zip code acts as a proxy for race and other protected
	attributes.
	B and D are incorrect. Temporal phenomena not indicated.
	C is incorrect. Overfitting is unrelated.
	
	\textbf{Learning Objective:} Describe sources of algorithmic bias.
	\textbf{Reference:} Module 3
	
	% =====================================================
	\subsection*{Question 19}
	Which fairness metric requires equal false positive rates across groups?
	
	A. Demographic parity, which equalizes positive prediction rates.\\
	B. Equalized odds, which equalizes both TPR and FPR across groups.\\
	C. Predictive parity, which equalizes positive predictive value.\\
	D. Calibration, which aligns predicted probabilities with outcomes.
	
	\vspace{0.4em}
	\textbf{Correct Answer: B}
	
	\textbf{Explanation:}
	A is incorrect. Demographic parity concerns positive prediction rates.
	B is correct. Equalized odds requires equal TPR and FPR.
	C and D are incorrect. Different fairness definitions.
	
	\textbf{Learning Objective:} Describe various measures of group fairness.
	\textbf{Reference:} Module 3
	
	% =====================================================
	\subsection*{Question 20}
	A regulator asks for the reasoning behind a credit denial. The bank
	cannot provide one. This most directly exposes the bank to:
	
	A. Market risk, arising from adverse price movements in positions.\\
	B. Liquidity risk, arising from an inability to meet cash obligations.\\
	C. Conduct and reputational risk, arising from unexplained decisions.\\
	D. Settlement risk, arising from failed counterparty transactions.
	
	\vspace{0.4em}
	\textbf{Correct Answer: C}
	
	\textbf{Explanation:}
	A, B, D are incorrect. Unrelated exposures.
	C is correct. Failure to explain decisions triggers conduct and
	reputational exposure.
	
	\textbf{Learning Objective:} Describe sources of AI-related reputational
	risk.
	\textbf{Reference:} Module 3
	
	% =====================================================
	\subsection*{Question 21}
	Which is the best approach to address historical bias in training data?
	
	A. Remove the protected attribute from the training dataset.\\
	B. Re-weight samples, audit outcomes, and document remediation.\\
	C. Increase the model's complexity to capture more patterns.\\
	D. Restrict training data to only the most recent observations.
	
	\vspace{0.4em}
	\textbf{Correct Answer: B}
	
	\textbf{Explanation:}
	A is incorrect. Removing the attribute does not remove the bias.
	B is correct. Re-weighting, auditing, and documenting is standard.
	C and D are incorrect. Not effective remediations.
	
	\textbf{Learning Objective:} Discuss sources of algorithmic bias.
	\textbf{Reference:} Module 4
	
	% =====================================================
	\subsection*{Question 22}
	A model shows high accuracy overall but poor recall for defaults. This
	is most concerning because:
	
	A. Accuracy is misleading in datasets with significant class imbalance.\\
	B. Precision is always more important than recall in credit models.\\
	C. Recall is only meaningful for continuous regression targets.\\
	D. Accuracy is a fairness metric that ignores class labels.
	
	\vspace{0.4em}
	\textbf{Correct Answer: A}
	
	\textbf{Explanation:}
	A is correct. High accuracy can mask poor minority-class recall.
	B is incorrect. Recall is often more important.
	C is incorrect. Recall is a classification metric.
	D is incorrect. Accuracy is not a fairness metric.
	
	\textbf{Learning Objective:} Evaluate classification model performance.
	\textbf{Reference:} Module 2, Chapter 8
	
	% =====================================================
	\subsection*{Question 23}
	An AI system used in hiring produces outcomes that disadvantage a
	protected group. Under most regulatory frameworks, the firm must:
	
	A. Take no action if overall accuracy is high for all candidates.\\
	B. Document, remediate, and demonstrate non-discriminatory outcomes.\\
	C. Publicly disclose the source code and training data for the model.\\
	D. Replace the model entirely with purely human hiring decisions.
	
	\vspace{0.4em}
	\textbf{Correct Answer: B}
	
	\textbf{Explanation:}
	A is incorrect. Accuracy does not excuse discrimination.
	B is correct. Documentation, remediation, and demonstration.
	C and D are incorrect. Not required.
	
	\textbf{Learning Objective:} Discuss sources of algorithmic bias.
	\textbf{Reference:} Module 4
	
	% =====================================================
	\subsection*{Question 24}
	Which best describes ``fairness through unawareness''?
	
	A. Removing protected attributes from the model's feature set.\\
	B. Training the model only on protected attributes and outcomes.\\
	C. Using demographic data as the sole basis for decisions.\\
	D. Ignoring all fairness metrics during model development.
	
	\vspace{0.4em}
	\textbf{Correct Answer: A}
	
	\textbf{Explanation:}
	A is correct. Fairness through unawareness = omitting protected
	attributes.
	B, C, D are incorrect. They are the opposite or irrelevant.
	
	\textbf{Learning Objective:} Describe various measures of group fairness.
	\textbf{Reference:} Module 3
	
	% =====================================================
	\subsection*{Question 25}
	The primary purpose of an AI ethics review board is to:
	
	A. Approve model deployment without requiring further review.\\
	B. Provide independent oversight of ethical implications.\\
	C. Approve marketing materials for AI-related products.\\
	D. Set interest rates for AI-driven credit products.
	
	\vspace{0.4em}
	\textbf{Correct Answer: B}
	
	\textbf{Explanation:}
	A is incorrect. The board reviews, does not rubber-stamp.
	B is correct. Independent ethical oversight.
	C and D are incorrect. Outside the board's remit.
	
	\textbf{Learning Objective:} Discuss the principles of AI ethics.
	\textbf{Reference:} Module 4
	
	% =====================================================
	\subsection*{Question 26}
	A model's performance differs across age groups. This is:
	
	A. Concept drift, caused by a change in the input-output relationship.\\
	B. Disparate impact, caused by differential outcomes across groups.\\
	C. Overfitting, caused by excessive model complexity and memorization.\\
	D. Underfitting, caused by insufficient model capacity to fit the data.
	
	\vspace{0.4em}
	\textbf{Correct Answer: B}
	
	\textbf{Explanation:}
	A is incorrect. Concept drift is temporal.
	B is correct. Different outcomes across groups = disparate impact.
	C and D are incorrect. Training issues, not group-specific.
	
	\textbf{Learning Objective:} Discuss sources of algorithmic bias.
	\textbf{Reference:} Module 3
	
	% =====================================================
	\subsection*{Question 27}
	Which best mitigates bias in an AI credit scoring system?
	
	A. Use only alternative data sources to replace traditional features.\\
	B. Combine fairness testing, diverse data, and human oversight.\\
	C. Use deeper neural networks to capture more complex patterns.\\
	D. Remove all features that correlate with protected attributes.
	
	\vspace{0.4em}
	\textbf{Correct Answer: B}
	
	\textbf{Explanation:}
	A is incorrect. Alternative data introduces its own bias.
	B is correct. Multi-layered mitigation.
	C is incorrect. Deep networks can amplify bias.
	D is incorrect. Removing features destroys utility.
	
	\textbf{Learning Objective:} Discuss sources of algorithmic bias.
	\textbf{Reference:} Module 4
	
	% =====================================================
	\subsection*{Question 28}
	Which is a human autonomy risk of AI in finance?
	
	A. Reduced transaction costs from automated processing.\\
	B. Decisions made without meaningful human review.\\
	C. Improved regulatory reporting through automation.\\
	D. Faster settlement of securities transactions.
	
	\vspace{0.4em}
	\textbf{Correct Answer: B}
	
	\textbf{Explanation:}
	A, C, D are incorrect. These are benefits, not autonomy risks.
	B is correct. Loss of meaningful review threatens autonomy.
	
	\textbf{Learning Objective:} Discuss risks posed by AI to human autonomy.
	\textbf{Reference:} Module 3
	
	% =====================================================
	\subsection*{Question 29}
	A chatbot gives medical advice beyond its intended scope. This most
	directly creates:
	
	A. Market risk, from adverse movements in traded positions.\\
	B. Conduct risk, from inappropriate customer interactions.\\
	C. Liquidity risk, from an inability to meet cash demands.\\
	D. Interest rate risk, from adverse rate movements.
	
	\vspace{0.4em}
	\textbf{Correct Answer: B}
	
	\textbf{Explanation:}
	A, C, D are incorrect. Unrelated risks.
	B is correct. Out-of-scope advice is a conduct issue.
	
	\textbf{Learning Objective:} Discuss risks posed by AI to human autonomy.
	\textbf{Reference:} Module 3
	
	% =====================================================
	\subsection*{Question 30}
	Which is the best way to increase trust in an AI system?
	
	A. Publish the model weights for full public inspection.\\
	B. Provide transparent explanations and clear audit trails.\\
	C. Increase model complexity to improve predictive accuracy.\\
	D. Reduce testing frequency to speed up deployment cycles.
	
	\vspace{0.4em}
	\textbf{Correct Answer: B}
	
	\textbf{Explanation:}
	A is incorrect. Weights alone do not build trust.
	B is correct. Transparency and audit trails build trust.
	C and D are incorrect. Reduce trust.
	
	\textbf{Learning Objective:} Describe explainability, interpretability,
	and transparency.
	\textbf{Reference:} Module 3
	
	% =====================================================
	\subsection*{Question 31}
	Which is a valid criticism of using accuracy alone to evaluate a
	fairness-sensitive model?
	
	A. Accuracy is not a legitimate performance metric in classification.\\
	B. Accuracy can hide poor performance for under-represented groups.\\
	C. Accuracy applies only to regression models, not classifiers.\\
	D. Accuracy ignores the class labels used in training the model.
	
	\vspace{0.4em}
	\textbf{Correct Answer: B}
	
	\textbf{Explanation:}
	A, C, D are incorrect. Factually wrong.
	B is correct. Accuracy can mask subgroup failures.
	
	\textbf{Learning Objective:} Evaluate classification model performance.
	\textbf{Reference:} Module 2, Chapter 8
	
	% =====================================================
	\subsection*{Question 32}
	Which AI risk is most associated with deepfakes in financial markets?
	
	A. Interest rate risk, from adverse movements in benchmark rates.\\
	B. Manipulation and misinformation risk, from fabricated content.\\
	C. Credit spread risk, from deterioration in issuer credit quality.\\
	D. Sovereign risk, from adverse government policy changes.
	
	\vspace{0.4em}
	\textbf{Correct Answer: B}
	
	\textbf{Explanation:}
	A, C, D are incorrect. Unrelated.
	B is correct. Deepfakes enable market manipulation and misinformation.
	
	\textbf{Learning Objective:} Discuss risks posed by AI to human autonomy.
	\textbf{Reference:} Module 3
	
	% =====================================================
	\subsection*{Question 33}
	A model's explanations are post-hoc approximations. This means:
	
	A. The explanations are mathematically exact and fully faithful.\\
	B. The explanations approximate the model's reasoning, not its true logic.\\
	C. The model produces no outputs that can be explained at all.\\
	D. The model is underfitting and therefore unreliable.
	
	\vspace{0.4em}
	\textbf{Correct Answer: B}
	
	\textbf{Explanation:}
	A is incorrect. Post-hoc explanations are approximate.
	B is correct. They approximate rather than reveal true logic.
	C and D are incorrect. Unrelated to post-hoc explanation.
	
	\textbf{Learning Objective:} Describe explainability, interpretability,
	and transparency.
	\textbf{Reference:} Module 3
	
	% =====================================================
	\subsection*{Question 34}
	Which framework best supports organizational AI ethics?
	
	A. Ad-hoc committees convened only when issues arise.\\
	B. Documented principles, review processes, and accountability.\\
	C. Outsourcing all ethical decisions to external vendors.\\
	D. Avoiding documentation to preserve proprietary methods.
	
	\vspace{0.4em}
	\textbf{Correct Answer: B}
	
	\textbf{Explanation:}
	A, C, D are incorrect. Not robust frameworks.
	B is correct. Documented principles + review + accountability.
	
	\textbf{Learning Objective:} Discuss the principles of AI ethics.
	\textbf{Reference:} Module 4
	
	% =====================================================
	\subsection*{Question 35}
	Which is the best action when a deployed AI produces unexpected
	outcomes for a subgroup?
	
	A. Ignore the issue if overall metrics remain strong.\\
	B. Escalate, investigate, and document remediation steps.\\
	C. Delete the affected logs to avoid regulatory exposure.\\
	D. Increase the model size to improve subgroup performance.
	
	\vspace{0.4em}
	\textbf{Correct Answer: B}
	
	\textbf{Explanation:}
	A is incorrect. Overall metrics can hide subgroup issues.
	B is correct. Escalate, investigate, document.
	C is incorrect. Deleting logs is a governance failure.
	D is incorrect. Model size is irrelevant.
	
	\textbf{Learning Objective:} Discuss sources of algorithmic bias.
	\textbf{Reference:} Module 4
	
	\newpage
	\section{Module 4: Ethics, Privacy, and Principles}
	
	\subsection*{Question 36}
	Which statement correctly describes Utilitarianism?
	
	A. It can lead to violations of individual rights for the greater good.\\
	B. It focuses on virtuous character traits and living a good life.\\
	C. It relies heavily on laws, regulations, and duty-based rules.\\
	D. It values the intrinsic morality of actions over their outcomes.
	
	\vspace{0.4em}
	\textbf{Correct Answer: A}
	
	\textbf{Explanation:}
	A is correct. Aggregate welfare can justify violating individual rights.
	B is incorrect. That describes virtue ethics.
	C is incorrect. That describes deontology.
	D is incorrect. That also describes deontology.
	
	\textbf{Learning Objective:} Compare ethical frameworks.
	\textbf{Reference:} Module 4
	
	% =====================================================
	\subsection*{Question 37}
	A deontological framework judges morality based on:
	
	A. The consequences of the action for total social welfare.\\
	B. Adherence to ethical duties and rules regardless of outcome.\\
	C. A single, quantifiable metric for overall utility.\\
	D. Virtuous character traits and the pursuit of a good life.
	
	\vspace{0.4em}
	\textbf{Correct Answer: B}
	
	\textbf{Explanation:}
	A is incorrect. Consequences = consequentialism.
	B is correct. Deontology = duty-based.
	C is incorrect. Utility metric = consequentialism.
	D is incorrect. Virtue = virtue ethics.
	
	\textbf{Learning Objective:} Compare ethical frameworks.
	\textbf{Reference:} Module 4
	
	% =====================================================
	\subsection*{Question 38}
	Regarding the principle of justice in AI ethics:
	
	A. Different concepts of distributive justice can conflict with one another.\\
	B. Impartial application of the law is optional for procedural justice.\\
	C. Personal autonomy is generally secondary to achieving social justice.\\
	D. Meritocracy is not considered a principle of distributive justice.
	
	\vspace{0.4em}
	\textbf{Correct Answer: A}
	
	\textbf{Explanation:}
	A is correct. Egalitarianism and utilitarianism can conflict.
	B is incorrect. Impartial application is required.
	C is incorrect. Social justice requires autonomy.
	D is incorrect. Meritocracy is a distributive principle.
	
	\textbf{Learning Objective:} Discuss the principles of AI ethics.
	\textbf{Reference:} Module 4
	
	% =====================================================
	\subsection*{Question 39}
	An AI recommends treatment for an unconscious patient that conflicts
	with documented preferences. Which principle is violated?
	
	A. Explainability, which concerns clarity of the AI's reasoning.\\
	B. Beneficence, which concerns acting in the patient's interest.\\
	C. Autonomy, which concerns respecting the patient's choices.\\
	D. Non-maleficence, which concerns avoiding harm to the patient.
	
	\vspace{0.4em}
	\textbf{Correct Answer: C}
	
	\textbf{Explanation:}
	A, B, D are incorrect. Secondary in this scenario.
	C is correct. Ignoring documented preferences violates autonomy.
	
	\textbf{Learning Objective:} Discuss the principles of AI ethics.
	\textbf{Reference:} Module 4
	
	% =====================================================
	\subsection*{Question 40}
	A model trained on data from the US and Canada performs poorly in
	Brazil, Nigeria, and Thailand. The most likely cause is:
	
	A. Bias in problem specification, from a mismatch in target variables.\\
	B. Bias in model deployment, from a failure in the rollout process.\\
	C. Historical bias, from outdated patterns in the training data.\\
	D. Sampling bias, from a non-representative training population.
	
	\vspace{0.4em}
	\textbf{Correct Answer: D}
	
	\textbf{Explanation:}
	A is incorrect. Not about target variables.
	B is incorrect. Not about deployment.
	C is incorrect. Not about time.
	D is correct. Non-representative training sample.
	
	\textbf{Learning Objective:} Discuss sources of algorithmic bias.
	\textbf{Reference:} Module 4
	
	% =====================================================
	\subsection*{Question 41}
	A model trained on 2005--2015 data underperforms in 2026 because
	societal patterns changed. This is:
	
	A. Sampling bias, from a mismatch between training and deployment populations.\\
	B. Historical bias, from using outdated data reflecting past conditions.\\
	C. Measurement bias, from inconsistent data collection methods over time.\\
	D. Data composition bias, from an imbalance in the training dataset.
	
	\vspace{0.4em}
	\textbf{Correct Answer: B}
	
	\textbf{Explanation:}
	A is incorrect. Not about geographic representation.
	B is correct. Outdated data = historical bias.
	C and D are incorrect. Not indicated.
	
	\textbf{Learning Objective:} Discuss sources of algorithmic bias.
	\textbf{Reference:} Module 4
	
	% =====================================================
	\subsection*{Question 42}
	Which ethical principle requires avoiding harm?
	
	A. Beneficence, which concerns acting for others' benefit.\\
	B. Non-maleficence, which concerns avoiding harm to others.\\
	C. Justice, which concerns fair distribution of benefits.\\
	D. Autonomy, which concerns respecting individual decisions.
	
	\vspace{0.4em}
	\textbf{Correct Answer: B}
	
	\textbf{Explanation:}
	A is incorrect. Beneficence = act for benefit.
	B is correct. Non-maleficence = do no harm.
	C and D are incorrect. Different principles.
	
	\textbf{Learning Objective:} Discuss the principles of AI ethics.
	\textbf{Reference:} Module 4
	
	% =====================================================
	\subsection*{Question 43}
	A company sells customer data without consent. Which principle is most
	directly violated?
	
	A. Explainability, which concerns clarity of decisions made by AI.\\
	B. Privacy, which concerns control over personal information.\\
	C. Fairness, which concerns equal treatment across groups.\\
	D. Transparency, which concerns disclosure of model design.
	
	\vspace{0.4em}
	\textbf{Correct Answer: B}
	
	\textbf{Explanation:}
	A, C, D are incorrect. Secondary or unrelated.
	B is correct. Selling data without consent is a privacy violation.
	
	\textbf{Learning Objective:} Describe principles of privacy.
	\textbf{Reference:} Module 4
	
	% =====================================================
	\subsection*{Question 44}
	Which best describes virtue ethics?
	
	A. An approach that focuses on consequences for overall welfare.\\
	B. An approach that focuses on duties and rules above outcomes.\\
	C. An approach that focuses on character traits and living well.\\
	D. An approach that focuses on maximizing total utility.
	
	\vspace{0.4em}
	\textbf{Correct Answer: C}
	
	\textbf{Explanation:}
	A is incorrect. Consequences = consequentialism.
	B is incorrect. Duties = deontology.
	C is correct. Virtue ethics = character.
	D is incorrect. Utility = consequentialism.
	
	\textbf{Learning Objective:} Compare ethical frameworks.
	\textbf{Reference:} Module 4
	
	% =====================================================
	\subsection*{Question 45}
	Which is the primary concern of beneficence in AI ethics?
	
	A. Avoiding harm to individuals and groups affected by AI systems.\\
	B. Acting in the interest of others and promoting their wellbeing.\\
	C. Preserving individual autonomy in decisions involving AI systems.\\
	D. Distributing benefits and burdens fairly across all groups.
	
	\vspace{0.4em}
	\textbf{Correct Answer: B}
	
	\textbf{Explanation:}
	A is incorrect. Avoiding harm = non-maleficence.
	B is correct. Beneficence = act for benefit.
	C is incorrect. Autonomy = self-determination.
	D is incorrect. Justice = fair distribution.
	
	\textbf{Learning Objective:} Discuss the principles of AI ethics.
	\textbf{Reference:} Module 4
	
	% =====================================================
	\subsection*{Question 46}
	Which best describes the principle of autonomy?
	
	A. Respecting individuals' right to make their own decisions.\\
	B. Maximizing aggregate utility across all affected parties.\\
	C. Following established rules regardless of consequences.\\
	D. Acting in a virtuous manner consistent with good character.
	
	\vspace{0.4em}
	\textbf{Correct Answer: A}
	
	\textbf{Explanation:}
	A is correct. Autonomy = self-determination.
	B is incorrect. Utility = consequentialism.
	C is incorrect. Rules = deontology.
	D is incorrect. Virtue = virtue ethics.
	
	\textbf{Learning Objective:} Discuss the principles of AI ethics.
	\textbf{Reference:} Module 4
	
	% =====================================================
	\subsection*{Question 47}
	A firm uses AI to make lending decisions without human review. Which
	principle is most directly at risk?
	
	A. Non-maleficence, which concerns avoiding harm to borrowers.\\
	B. Autonomy, which concerns respecting borrowers' decision-making.\\
	C. Beneficence, which concerns acting in borrowers' interest.\\
	D. Justice, which concerns fair treatment of all borrowers.
	
	\vspace{0.4em}
	\textbf{Correct Answer: B}
	
	\textbf{Explanation:}
	A, C, D are incorrect. Secondary.
	B is correct. No human review threatens applicant autonomy.
	
	\textbf{Learning Objective:} Discuss the principles of AI ethics.
	\textbf{Reference:} Module 4
	
	% =====================================================
	\subsection*{Question 48}
	Which best describes informational privacy?
	
	A. Control over how personal data is collected and used.\\
	B. The right to physical security and bodily integrity.\\
	C. Ownership of intellectual property created by individuals.\\
	D. Protection of freedom of speech and expression.
	
	\vspace{0.4em}
	\textbf{Correct Answer: A}
	
	\textbf{Explanation:}
	A is correct. Informational privacy = control over personal data.
	B, C, D are incorrect. Different concepts.
	
	\textbf{Learning Objective:} Describe principles of privacy.
	\textbf{Reference:} Module 4
	
	% =====================================================
	\subsection*{Question 49}
	Which is a legitimate basis for processing personal data under most
	privacy frameworks?
	
	A. Any business interest that the firm deems important.\\
	B. Explicit consent or a legal obligation to process.\\
	C. Speculation about future uses of the data.\\
	D. A general agreement with a third-party vendor.
	
	\vspace{0.4em}
	\textbf{Correct Answer: B}
	
	\textbf{Explanation:}
	A, C, D are incorrect. Not sufficient legal bases.
	B is correct. Consent or legal obligation.
	
	\textbf{Learning Objective:} Describe principles of privacy.
	\textbf{Reference:} Module 4
	
	% =====================================================
	\subsection*{Question 50}
	Which best describes data minimization?
	
	A. Collecting as much data as possible to enable future use.\\
	B. Collecting only data necessary for the stated purpose.\\
	C. Storing data indefinitely to preserve analytical value.\\
	D. Sharing data with all vendors that support the business.
	
	\vspace{0.4em}
	\textbf{Correct Answer: B}
	
	\textbf{Explanation:}
	A, C, D are incorrect. Contradict minimization.
	B is correct. Minimization = collect only what is needed.
	
	\textbf{Learning Objective:} Describe principles of privacy.
	\textbf{Reference:} Module 4
	
	% =====================================================
	\subsection*{Question 51}
	A firm retains customer data longer than necessary. This violates:
	
	A. Data minimization, which limits collection and retention.\\
	B. Autonomy, which concerns respecting individual decisions.\\
	C. Beneficence, which concerns acting in others' interest.\\
	D. Transparency, which concerns disclosure of data practices.
	
	\vspace{0.4em}
	\textbf{Correct Answer: A}
	
	\textbf{Explanation:}
	A is correct. Retention beyond necessity violates minimization.
	B, C, D are incorrect. Secondary or unrelated.
	
	\textbf{Learning Objective:} Describe principles of privacy.
	\textbf{Reference:} Module 4
	
	% =====================================================
	\subsection*{Question 52}
	Which best describes deontology?
	
	A. An ethical framework based on the consequences of actions.\\
	B. An ethical framework based on duties and moral rules.\\
	C. An ethical framework based on character and virtue.\\
	D. An ethical framework based on maximizing total utility.
	
	\vspace{0.4em}
	\textbf{Correct Answer: B}
	
	\textbf{Explanation:}
	A is incorrect. Consequence-based = consequentialism.
	B is correct. Duty-based = deontology.
	C is incorrect. Character = virtue ethics.
	D is incorrect. Utility = consequentialism.
	
	\textbf{Learning Objective:} Compare ethical frameworks.
	\textbf{Reference:} Module 4
	
	% =====================================================
	\subsection*{Question 53}
	Which best describes consequentialism?
	
	A. Actions are judged by their outcomes and overall effects.\\
	B. Actions are judged by adherence to duties and rules.\\
	C. Actions are judged by the character of the actor.\\
	D. Actions are judged by their conformity to laws.
	
	\vspace{0.4em}
	\textbf{Correct Answer: A}
	
	\textbf{Explanation:}
	A is correct. Consequentialism = outcomes.
	B is incorrect. Duties = deontology.
	C is incorrect. Character = virtue ethics.
	D is incorrect. Laws = legal positivism.
	
	\textbf{Learning Objective:} Compare ethical frameworks.
	\textbf{Reference:} Module 4
	
	% =====================================================
	\subsection*{Question 54}
	A firm uses AI to determine credit limits based on social media
	activity. The key ethical concern is:
	
	A. Privacy and consent, from using non-financial personal data.\\
	B. Interest rate levels, from adverse rate movements.\\
	C. Model inference speed, from computational constraints.\\
	D. Data storage cost, from increased storage requirements.
	
	\vspace{0.4em}
	\textbf{Correct Answer: A}
	
	\textbf{Explanation:}
	A is correct. Social media data raises consent/privacy concerns.
	B, C, D are incorrect. Unrelated to the ethics issue.
	
	\textbf{Learning Objective:} Describe principles of privacy.
	\textbf{Reference:} Module 4
	
	% =====================================================
	\subsection*{Question 55}
	Which is a risk of using alternative data in credit models?
	
	A. Only improved inclusion of previously unscorable applicants.\\
	B. Potential use of sensitive or non-consented personal data.\\
	C. Reduced model accuracy relative to traditional data.\\
	D. Lower regulatory scrutiny of the credit decision process.
	
	\vspace{0.4em}
	\textbf{Correct Answer: B}
	
	\textbf{Explanation:}
	A is incorrect. Inclusion is a potential benefit, not the only effect.
	B is correct. Alternative data may be sensitive or non-consented.
	C is incorrect. Accuracy may improve.
	D is incorrect. More scrutiny, not less.
	
	\textbf{Learning Objective:} Describe data governance frameworks.
	\textbf{Reference:} Module 5
	
	% =====================================================
	\subsection*{Question 56}
	Which best describes the principle of transparency?
	
	A. Publishing model weights and full training code publicly.\\
	B. Providing accessible information about how decisions are made.\\
	C. Sharing all training data with affected individuals.\\
	D. Publishing source code without documentation.
	
	\vspace{0.4em}
	\textbf{Correct Answer: B}
	
	\textbf{Explanation:}
	A, C, D are incorrect. Specific actions that may not constitute
	transparency.
	B is correct. Transparency = accessible information on decisions.
	
	\textbf{Learning Objective:} Describe explainability, interpretability,
	and transparency.
	\textbf{Reference:} Module 3
	
	% =====================================================
	\subsection*{Question 57}
	Which is the best approach when ethical principles conflict?
	
	A. Ignore the conflict and proceed with the fastest option.\\
	B. Apply a documented prioritization and governance framework.\\
	C. Follow the option that maximizes short-term profit.\\
	D. Choose arbitrarily to avoid delaying deployment.
	
	\vspace{0.4em}
	\textbf{Correct Answer: B}
	
	\textbf{Explanation:}
	A, C, D are incorrect. Not ethical responses.
	B is correct. Documented prioritization and governance.
	
	\textbf{Learning Objective:} Discuss the principles of AI ethics.
	\textbf{Reference:} Module 4
	
	% =====================================================
	\subsection*{Question 58}
	Which best describes the role of consent in AI?
	
	A. Consent is optional in all cases involving personal data.\\
	B. Consent is required when no other legal basis exists.\\
	C. Consent is never required for commercial AI applications.\\
	D. Consent is only required for government AI applications.
	
	\vspace{0.4em}
	\textbf{Correct Answer: B}
	
	\textbf{Explanation:}
	A, C, D are incorrect. Consent is generally required absent legal basis.
	B is correct. Consent required if no other legal basis.
	
	\textbf{Learning Objective:} Describe principles of privacy.
	\textbf{Reference:} Module 4
	
	% =====================================================
	\subsection*{Question 59}
	Which best describes algorithmic fairness?
	
	A. Ensuring identical outcomes for every individual in a population.\\
	B. Ensuring fair treatment across groups based on chosen criteria.\\
	C. Ensuring maximum accuracy regardless of subgroup performance.\\
	D. Ensuring no automated model is used in high-stakes decisions.
	
	\vspace{0.4em}
	\textbf{Correct Answer: B}
	
	\textbf{Explanation:}
	A is incorrect. Identical outcomes are not always the goal.
	B is correct. Fair treatment based on chosen fairness criteria.
	C is incorrect. Accuracy alone is not fairness.
	D is incorrect. Not a definition of fairness.
	
	\textbf{Learning Objective:} Describe various measures of group fairness.
	\textbf{Reference:} Module 3
	
	% =====================================================
	\subsection*{Question 60}
	Which best describes non-maleficence?
	
	A. Do no harm to individuals and groups affected by AI.\\
	B. Do good by acting in the interest of others.\\
	C. Act fairly by distributing benefits equitably.\\
	D. Respect autonomy by honoring individual choices.
	
	\vspace{0.4em}
	\textbf{Correct Answer: A}
	
	\textbf{Explanation:}
	A is correct. Non-maleficence = do no harm.
	B is incorrect. Beneficence = do good.
	C is incorrect. Justice = act fairly.
	D is incorrect. Autonomy = respect choice.
	
	\textbf{Learning Objective:} Discuss the principles of AI ethics.
	\textbf{Reference:} Module 4
	
	% =====================================================
	\subsection*{Question 61}
	Which best describes beneficence?
	
	A. Acting in a way that benefits others and promotes wellbeing.\\
	B. Avoiding harm to individuals and groups affected by AI.\\
	C. Respecting individuals' rights to make their own decisions.\\
	D. Distributing benefits and burdens fairly across all groups.
	
	\vspace{0.4em}
	\textbf{Correct Answer: A}
	
	\textbf{Explanation:}
	A is correct. Beneficence = act for benefit.
	B is incorrect. Avoid harm = non-maleficence.
	C is incorrect. Respect rights = autonomy.
	D is incorrect. Fair distribution = justice.
	
	\textbf{Learning Objective:} Discuss the principles of AI ethics.
	\textbf{Reference:} Module 4
	
	% =====================================================
	\subsection*{Question 62}
	A firm uses AI to make employment decisions but keeps the algorithm
	secret. This most directly affects:
	
	A. Transparency, which concerns accessible information about decisions.\\
	B. Autonomy, which concerns respecting individual decision-making.\\
	C. Beneficence, which concerns acting in others' interest.\\
	D. Non-maleficence, which concerns avoiding harm to others.
	
	\vspace{0.4em}
	\textbf{Correct Answer: A}
	
	\textbf{Explanation:}
	A is correct. Secrecy undermines transparency.
	B, C, D are incorrect. Secondary.
	
	\textbf{Learning Objective:} Describe explainability, interpretability,
	and transparency.
	\textbf{Reference:} Module 3
	
	% =====================================================
	\subsection*{Question 63}
	Which best describes distributive justice?
	
	A. Fair allocation of benefits and burdens across a society.\\
	B. Fair application of rules and procedures to all.\\
	C. Fair development of individual character and virtue.\\
	D. Fair maximization of aggregate utility and welfare.
	
	\vspace{0.4em}
	\textbf{Correct Answer: A}
	
	\textbf{Explanation:}
	A is correct. Distributive justice = fair allocation.
	B is incorrect. Fair rules = procedural.
	C is incorrect. Character = virtue ethics.
	D is incorrect. Utility = consequentialism.
	
	\textbf{Learning Objective:} Discuss the principles of justice.
	\textbf{Reference:} Module 4
	
	% =====================================================
	\subsection*{Question 64}
	Which best describes retributive justice?
	
	A. Fair punishment for wrongdoing proportional to the offense.\\
	B. Fair allocation of resources across all societal groups.\\
	C. Fair procedural application of rules and laws.\\
	D. Fair development of virtuous character in individuals.
	
	\vspace{0.4em}
	\textbf{Correct Answer: A}
	
	\textbf{Explanation:}
	A is correct. Retributive = just punishment.
	B is incorrect. Allocation = distributive.
	C is incorrect. Procedure = procedural.
	D is incorrect. Character = virtue ethics.
	
	\textbf{Learning Objective:} Discuss the principles of justice.
	\textbf{Reference:} Module 4
	
	% =====================================================
	\subsection*{Question 65}
	Which best describes procedural justice?
	
	A. Fair processes for making and enforcing decisions.\\
	B. Fair outcomes in the distribution of benefits.\\
	C. Fair development of individual character traits.\\
	D. Fair maximization of overall social welfare.
	
	\vspace{0.4em}
	\textbf{Correct Answer: A}
	
	\textbf{Explanation:}
	A is correct. Procedural justice = fair processes.
	B is incorrect. Outcomes = distributive.
	C, D are incorrect. Unrelated.
	
	\textbf{Learning Objective:} Discuss the principles of justice.
	\textbf{Reference:} Module 4
	
	% =====================================================
	\subsection*{Question 66}
	Which is the best response to a privacy breach caused by AI?
	
	A. Ignore the breach if the number of affected users is small.\\
	B. Notify affected parties, remediate, and document the incident.\\
	C. Delete the affected logs to prevent further disclosure.\\
	D. Publicly blame users for sharing their personal information.
	
	\vspace{0.4em}
	\textbf{Correct Answer: B}
	
	\textbf{Explanation:}
	A, C, D are incorrect. Irresponsible and potentially illegal.
	B is correct. Notification + remediation + documentation.
	
	\textbf{Learning Objective:} Describe principles of privacy.
	\textbf{Reference:} Module 4
	
	% =====================================================
	\subsection*{Question 67}
	Which best describes data subject rights?
	
	A. The rights of individuals over their own personal data.\\
	B. The rights of firms over customer data they collect.\\
	C. The rights of regulators to access corporate data.\\
	D. The rights of vendors to use data shared by clients.
	
	\vspace{0.4em}
	\textbf{Correct Answer: A}
	
	\textbf{Explanation:}
	A is correct. Data subject = the individual.
	B, C, D are incorrect. Not the standard definition.
	
	\textbf{Learning Objective:} Describe principles of privacy.
	\textbf{Reference:} Module 4
	
	% =====================================================
	\subsection*{Question 68}
	Which is a potential harm of AI in hiring?
	
	A. Reduced bias in hiring decisions through automation.\\
	B. Amplified historical discrimination against certain groups.\\
	C. Increased salaries for all candidates who apply.\\
	D. Lower recruiting costs for the hiring organization.
	
	\vspace{0.4em}
	\textbf{Correct Answer: B}
	
	\textbf{Explanation:}
	A is incorrect. Bias is not automatically reduced.
	B is correct. Historical discrimination can be amplified.
	C, D are incorrect. These are not harms.
	
	\textbf{Learning Objective:} Discuss sources of algorithmic bias.
	\textbf{Reference:} Module 4
	
	% =====================================================
	\subsection*{Question 69}
	Which best describes the principle of accountability in AI?
	
	A. No single party is responsible for AI outcomes.\\
	B. Clear responsibility for AI outcomes and their consequences.\\
	C. Only the developers are responsible for AI outcomes.\\
	D. Only the end users are responsible for AI outcomes.
	
	\vspace{0.4em}
	\textbf{Correct Answer: B}
	
	\textbf{Explanation:}
	A, C, D are incorrect. Incomplete or wrong.
	B is correct. Accountability = clear responsibility.
	
	\textbf{Learning Objective:} Discuss the principles of AI ethics.
	\textbf{Reference:} Module 4
	
	% =====================================================
	\subsection*{Question 70}
	Which best describes the ethical concern with autonomous AI
	decision-making?
	
	A. Loss of human oversight and clear accountability for decisions.\\
	B. Increased transparency and explainability of decisions.\\
	C. Reduced operating cost from automation of decisions.\\
	D. Improved fairness from consistent algorithmic treatment.
	
	\vspace{0.4em}
	\textbf{Correct Answer: A}
	
	\textbf{Explanation:}
	A is correct. Loss of oversight and accountability is the concern.
	B, C, D are incorrect. These are benefits, not concerns.
	
	\textbf{Learning Objective:} Discuss risks posed by AI to human autonomy.
	\textbf{Reference:} Module 3
	
	\newpage
	\section{Module 5: Model Risk Management and Governance}
	
	\subsection*{Question 71}
	A local bank periodically validates its VaR model. Which is a best
	practice?
	
	A. Validation should always occur on a fixed quarterly schedule.\\
	B. Validation should assess whether original assumptions remain valid.\\
	C. Validation is discretionary and not required by regulators.\\
	D. Validation should be conducted by the bank's CRO personally.
	
	\vspace{0.4em}
	\textbf{Correct Answer: B}
	
	\textbf{Explanation:}
	A is incorrect. Frequency depends on model risk ranking.
	B is correct. Validation checks assumptions, data, and performance.
	C is incorrect. Validation is regulatorily prescribed.
	D is incorrect. No such requirement.
	
	\textbf{Learning Objective:} Discuss model validation and its importance.
	\textbf{Reference:} Module 5
	
	% =====================================================
	\subsection*{Question 72}
	A CRO evaluates the model landscape. The most crucial consideration is:
	
	A. Ensuring that each model operates fully independently of others.\\
	B. Ensuring that all models use the same set of input data.\\
	C. Ensuring that models can be consolidated to reduce complexity.\\
	D. Ensuring that interdependencies between models are identified.
	
	\vspace{0.4em}
	\textbf{Correct Answer: D}
	
	\textbf{Explanation:}
	A is incorrect. Independence is not required.
	B is incorrect. Same inputs are not required.
	C is incorrect. Consolidation is not the focus.
	D is correct. Interdependencies are the focus of landscape review.
	
	\textbf{Learning Objective:} Describe factors for model inventory.
	\textbf{Reference:} Module 5
	
	% =====================================================
	\subsection*{Question 73}
	Which correctly describes appropriate roles in model governance?
	
	A. The board of directors should maintain the model inventory.\\
	B. The model risk function should set the model risk appetite.\\
	C. Model ownership should never be shared between individuals.\\
	D. Internal and external auditors are necessary independent checks.
	
	\vspace{0.4em}
	\textbf{Correct Answer: D}
	
	\textbf{Explanation:}
	A is incorrect. Model risk function maintains the inventory.
	B is incorrect. The board sets model risk appetite.
	C is incorrect. Ownership may be shared.
	D is correct. Auditors provide independent checks.
	
	\textbf{Learning Objective:} Discuss roles and responsibilities in MRM.
	\textbf{Reference:} Module 5
	
	% =====================================================
	\subsection*{Question 74}
	Which team ensures a new ML model integrates with IT systems without
	disrupting operations?
	
	A. The Core Implementation Team, which handles model algorithms.\\
	B. The System Implementation Team, which handles system integration.\\
	C. The Model Design Team, which handles model architecture.\\
	D. The Data Management Team, which handles data pipelines.
	
	\vspace{0.4em}
	\textbf{Correct Answer: B}
	
	\textbf{Explanation:}
	A, C, D are incorrect. Not their primary responsibility.
	B is correct. System Implementation Team integrates the model.
	
	\textbf{Learning Objective:} Describe steps in model implementation.
	\textbf{Reference:} Module 5
	
	% =====================================================
	\subsection*{Question 75}
	What is an advantage of increasing model monitoring frequency in rapidly
	changing environments?
	
	A. It reduces the need for model retraining over time.\\
	B. It detects performance degradation or drift for timely intervention.\\
	C. It minimizes the need for IT involvement in model operations.\\
	D. It improves the model's interpretability to business users.
	
	\vspace{0.4em}
	\textbf{Correct Answer: B}
	
	\textbf{Explanation:}
	A is incorrect. Retraining is still required.
	B is correct. Monitoring detects drift early.
	C is incorrect. Monitoring may increase IT involvement.
	D is incorrect. Monitoring does not change interpretability.
	
	\textbf{Learning Objective:} Describe AI/ML model review framework.
	\textbf{Reference:} Module 5
	
	% =====================================================
	\subsection*{Question 76}
	Relative to traditional statistical models, data quality requirements
	for AI/ML models are:
	
	A. Identical, because both require the same data standards.\\
	B. Less rigorous, because AI/ML models are more robust to noise.\\
	C. Focused only on outlier detection and removal procedures.\\
	D. More stringent, because AI/ML models use alternative data sources.
	
	\vspace{0.4em}
	\textbf{Correct Answer: D}
	
	\textbf{Explanation:}
	A, B, C are incorrect. AI/ML models have higher requirements.
	D is correct. AI/ML data quality requirements are more stringent.
	
	\textbf{Learning Objective:} Describe data governance frameworks.
	\textbf{Reference:} Module 5
	
	% =====================================================
	\subsection*{Question 77}
	A new AI mortgage prepayment model outperforms the existing statistical
	model in testing. The appropriate next step is:
	
	A. Deploy the model directly to production for immediate use.\\
	B. Notify the IT department to roll out the model to users.\\
	C. Retire the existing model and replace it with the new one.\\
	D. Request initial validation from the model validation team.
	
	\vspace{0.4em}
	\textbf{Correct Answer: D}
	
	\textbf{Explanation:}
	A, C are incorrect. Deployment requires validation first.
	B is incorrect. IT notification belongs to the design phase.
	D is correct. Initial validation is required before production.
	
	\textbf{Learning Objective:} Discuss model validation and its importance.
	\textbf{Reference:} Module 5, Section 3.2.1
	
	% =====================================================
	\subsection*{Question 78}
	A new AI credit model uses school transcripts, part-time jobs, social
	media, and landlord-tenant disputes. A validation finding should be:
	
	A. No material concerns, because augmented data improves inclusion.\\
	B. Material concerns about sensitive or non-consented alternative data.\\
	C. Approval, because backtesting shows improvement over the old model.\\
	D. Recommendation to include additional alternative data sources.
	
	\vspace{0.4em}
	\textbf{Correct Answer: B}
	
	\textbf{Explanation:}
	A is incorrect. Sensitive data raises concerns.
	B is correct. Sensitive alternative data should trigger material concerns.
	C is incorrect. Backtesting does not override data ethics concerns.
	D is incorrect. More alternative data increases risk.
	
	\textbf{Learning Objective:} Describe data governance frameworks.
	\textbf{Reference:} Module 5
	
	% =====================================================
	\subsection*{Question 79}
	Which is a typical element of a data governance framework for AI?
	
	A. Data minimization as the sole governing principle.\\
	B. Provenance, quality, access controls, and lifecycle management.\\
	C. Vendor contract terms as the primary governance mechanism.\\
	D. Cloud storage cost controls and usage optimization.
	
	\vspace{0.4em}
	\textbf{Correct Answer: B}
	
	\textbf{Explanation:}
	A, C, D are incorrect. Incomplete elements.
	B is correct. Provenance, quality, access, and lifecycle are core.
	
	\textbf{Learning Objective:} Describe data governance frameworks.
	\textbf{Reference:} Module 5
	
	% =====================================================
	\subsection*{Question 80}
	Model owners in complex organizations:
	
	A. Must always be a single individual to ensure accountability.\\
	B. May be shared among individuals or assigned to a committee.\\
	C. Cannot be assigned to committees under any circumstances.\\
	D. Are always external parties independent of the organization.
	
	\vspace{0.4em}
	\textbf{Correct Answer: B}
	
	\textbf{Explanation:}
	A, C, D are incorrect. Ownership may be shared.
	B is correct. Shared ownership is allowed in complex organizations.
	
	\textbf{Learning Objective:} Discuss roles and responsibilities in MRM.
	\textbf{Reference:} Module 5
	
	% =====================================================
	\subsection*{Question 81}
	Which is not a typical step in AI/ML model implementation?
	
	A. Model design, including architecture and feature selection.\\
	B. System integration with existing IT infrastructure.\\
	C. Ongoing monitoring after deployment into production.\\
	D. Permanent exemption from validation once deployed.
	
	\vspace{0.4em}
	\textbf{Correct Answer: D}
	
	\textbf{Explanation:}
	A, B, C are incorrect. These are typical steps.
	D is correct. Permanent exemption from validation is not a step.
	
	\textbf{Learning Objective:} Describe steps in model implementation.
	\textbf{Reference:} Module 5
	
	% =====================================================
	\subsection*{Question 82}
	Model risk appetite for a bank should be set by:
	
	A. The model risk function, based on technical analysis.\\
	B. The board of directors, as part of overall governance.\\
	C. The CRO alone, without board involvement.\\
	D. Internal audit, as an independent function.
	
	\vspace{0.4em}
	\textbf{Correct Answer: B}
	
	\textbf{Explanation:}
	A, C, D are incorrect. Not the correct authority.
	B is correct. The board sets model risk appetite.
	
	\textbf{Learning Objective:} Discuss roles and responsibilities in MRM.
	\textbf{Reference:} Module 5
	
	% =====================================================
	\subsection*{Question 83}
	A model inventory should be maintained by:
	
	A. The board of directors as part of strategic oversight.\\
	B. The model risk function as part of its mandate.\\
	C. Internal audit as part of its independent review.\\
	D. IT operations as part of system administration.
	
	\vspace{0.4em}
	\textbf{Correct Answer: B}
	
	\textbf{Explanation:}
	A, C, D are incorrect. Not the maintaining function.
	B is correct. Model risk function maintains the inventory.
	
	\textbf{Learning Objective:} Discuss roles and responsibilities in MRM.
	\textbf{Reference:} Module 5
	
	% =====================================================
	\subsection*{Question 84}
	The primary objective of periodic model validation is to:
	
	A. Confirm that the model continues to generate profit for the firm.\\
	B. Ensure assumptions remain valid and performance is as expected.\\
	C. Replace the model with a newer version on an annual cycle.\\
	D. Satisfy internal audit requirements for documentation.
	
	\vspace{0.4em}
	\textbf{Correct Answer: B}
	
	\textbf{Explanation:}
	A, C, D are incorrect. Not the objective.
	B is correct. Validation ensures assumptions and performance.
	
	\textbf{Learning Objective:} Discuss model validation and its importance.
	\textbf{Reference:} Module 5
	
	% =====================================================
	\subsection*{Question 85}
	Which is a risk unique to generative AI in risk management workflows?
	
	A. Interest rate exposure from adverse movements in rates.\\
	B. Hallucination producing fabricated facts and citations.\\
	C. Liquidity mismatch from timing differences in cash flows.\\
	D. Settlement risk from failed counterparty transactions.
	
	\vspace{0.4em}
	\textbf{Correct Answer: B}
	
	\textbf{Explanation:}
	A, C, D are incorrect. Not unique to GenAI.
	B is correct. Hallucination is a GenAI-specific risk.
	
	\textbf{Learning Objective:} Describe AI/ML model review framework.
	\textbf{Reference:} Module 5
	
	% =====================================================
	\subsection*{Question 86}
	Under the EU AI Act, credit scoring of individuals is classified as:
	
	A. Prohibited, because it involves personal financial data.\\
	B. High-risk, because it affects access to essential services.\\
	C. Limited-risk, because it only affects individual consumers.\\
	D. Minimal-risk, because it is a routine banking operation.
	
	\vspace{0.4em}
	\textbf{Correct Answer: B}
	
	\textbf{Explanation:}
	A is incorrect. Not prohibited.
	B is correct. Credit scoring is high-risk.
	C and D are incorrect. Too low.
	
	\textbf{Learning Objective:} Describe regulatory frameworks for AI.
	\textbf{Reference:} Module 5
	
	% =====================================================
	\subsection*{Question 87}
	A bank uses an LLM to summarize counterparty news for credit decisions.
	The most significant risk is:
	
	A. Hallucinated facts influencing the resulting credit decision.\\
	B. High electricity consumption from running the model.\\
	C. Volatility in token pricing affecting operational cost.\\
	D. The model's training data cut-off date being outdated.
	
	\vspace{0.4em}
	\textbf{Correct Answer: A}
	
	\textbf{Explanation:}
	A is correct. Hallucination is the dominant risk.
	B, C, D are incorrect. Secondary operational concerns.
	
	\textbf{Learning Objective:} Describe AI/ML model review framework.
	\textbf{Reference:} Module 5
	
	% =====================================================
	\subsection*{Question 88}
	Third-party AI risk primarily concerns:
	
	A. Vendor contract cost and pricing renegotiation exposure.\\
	B. Concentration, data leakage, model opacity, and dependency.\\
	C. Vendor equity price volatility affecting partnership value.\\
	D. Vendor exposure to changes in benchmark interest rates.
	
	\vspace{0.4em}
	\textbf{Correct Answer: B}
	
	\textbf{Explanation:}
	A, C, D are incorrect. Too narrow.
	B is correct. Multidimensional third-party risk.
	
	\textbf{Learning Objective:} Discuss vendor and third-party risk.
	\textbf{Reference:} Module 5
	
	% =====================================================
	\subsection*{Question 89}
	A model's predictions become less accurate over time because input
	distributions shift. This is:
	
	A. Concept drift, caused by changes in the input-output relationship.\\
	B. Data drift, caused by changes in the distribution of inputs.\\
	C. Overfitting, caused by excessive complexity and memorization.\\
	D. Underfitting, caused by insufficient model capacity.
	
	\vspace{0.4em}
	\textbf{Correct Answer: B}
	
	\textbf{Explanation:}
	A is incorrect. Concept drift = change in $P(y|X)$.
	B is correct. Data drift = change in $P(X)$.
	C, D are incorrect. Training issues.
	
	\textbf{Learning Objective:} Describe AI/ML model review framework.
	\textbf{Reference:} Module 5
	
	% =====================================================
	\subsection*{Question 90}
	Which best describes the relationship between interpretability and
	explainability?
	
	A. They are essentially identical concepts used interchangeably.\\
	B. Interpretability concerns model structure; explainability concerns
	individual predictions.\\
	C. Explainability applies only to linear models; interpretability is
	universal.\\
	D. Interpretability requires deep networks; explainability requires
	shallow models.
	
	\vspace{0.4em}
	\textbf{Correct Answer: B}
	
	\textbf{Explanation:}
	A is incorrect. Related but distinct.
	B is correct. Interpretability = model structure; explainability =
	individual predictions.
	C, D are incorrect. Reverse the relationship.
	
	\textbf{Learning Objective:} Describe explainability, interpretability,
	and transparency.
	\textbf{Reference:} Module 3
	
	% =====================================================
	\subsection*{Question 91}
	Which is the most appropriate governance response when a model's
	performance degrades significantly?
	
	A. Continue using it while documenting performance in quarterly reports.\\
	B. Trigger review, investigate root cause, and remediate or retire.\\
	C. Increase marketing efforts to compensate for lower performance.\\
	D. Ignore the degradation if the model was previously validated.
	
	\vspace{0.4em}
	\textbf{Correct Answer: B}
	
	\textbf{Explanation:}
	A, C, D are incorrect. Irresponsible.
	B is correct. Trigger review, investigate, remediate or retire.
	
	\textbf{Learning Objective:} Describe steps in model adaptation.
	\textbf{Reference:} Module 5
	
	% =====================================================
	\subsection*{Question 92}
	Which best describes the three lines of defense in model risk
	management?
	
	A. Business, model risk, internal audit.\\
	B. Marketing, sales, and human resources.\\
	C. IT, HR, and finance functions.\\
	D. Legal, compliance, and marketing.
	
	\vspace{0.4em}
	\textbf{Correct Answer: A}
	
	\textbf{Explanation:}
	A is correct. Business = first line; model risk = second; audit = third.
	B, C, D are incorrect. Not the three lines.
	
	\textbf{Learning Objective:} Discuss roles and responsibilities in MRM.
	\textbf{Reference:} Module 5
	
	% =====================================================
	\subsection*{Question 93}
	Which is the primary responsibility of the first line of defense in
	model risk?
	
	A. Independent validation of models before production use.\\
	B. Owning and managing model risk in day-to-day operations.\\
	C. Auditing the model governance framework independently.\\
	D. Setting regulatory policy for model risk management.
	
	\vspace{0.4em}
	\textbf{Correct Answer: B}
	
	\textbf{Explanation:}
	A is incorrect. Validation = second line.
	B is correct. First line owns and manages model risk.
	C is incorrect. Audit = third line.
	D is incorrect. Policy setting is not first line's role.
	
	\textbf{Learning Objective:} Discuss roles and responsibilities in MRM.
	\textbf{Reference:} Module 5
	
	% =====================================================
	\subsection*{Question 94}
	Which is the primary responsibility of the second line of defense?
	
	A. Model development and feature engineering.\\
	B. Independent model validation and oversight.\\
	C. Internal audit of the governance framework.\\
	D. Marketing of AI-based products.
	
	\vspace{0.4em}
	\textbf{Correct Answer: B}
	
	\textbf{Explanation:}
	A is incorrect. Development = first line.
	B is correct. Second line = independent validation and oversight.
	C is incorrect. Audit = third line.
	D is incorrect. Marketing is unrelated.
	
	\textbf{Learning Objective:} Discuss roles and responsibilities in MRM.
	\textbf{Reference:} Module 5
	
	% =====================================================
	\subsection*{Question 95}
	Which is the primary responsibility of the third line of defense?
	
	A. Model development and calibration activities.\\
	B. Independent audit of the model risk framework.\\
	C. Model validation and performance testing.\\
	D. Business use and daily operation of models.
	
	\vspace{0.4em}
	\textbf{Correct Answer: B}
	
	\textbf{Explanation:}
	A is incorrect. Development = first line.
	B is correct. Third line = independent audit.
	C is incorrect. Validation = second line.
	D is incorrect. Business use = first line.
	
	\textbf{Learning Objective:} Discuss roles and responsibilities in MRM.
	\textbf{Reference:} Module 5
	
	% =====================================================
	\subsection*{Question 96}
	A model's documentation is incomplete. This is a violation of:
	
	A. Model risk management standards requiring documentation.\\
	B. Accounting standards requiring financial disclosures.\\
	C. Marketing standards requiring consumer disclosures.\\
	D. Human resources standards requiring role clarity.
	
	\vspace{0.4em}
	\textbf{Correct Answer: A}
	
	\textbf{Explanation:}
	A is correct. Documentation is required by MRM standards.
	B, C, D are incorrect. Unrelated standards.
	
	\textbf{Learning Objective:} Discuss model validation and its importance.
	\textbf{Reference:} Module 5
	
	% =====================================================
	\subsection*{Question 97}
	Which is the best approach for validating an AI/ML model?
	
	A. Only statistical tests on in-sample prediction accuracy.\\
	B. Conceptual soundness, data quality, outcomes analysis, and monitoring.\\
	C. Only backtesting against historical performance data.\\
	D. Only vendor-provided reports and certifications.
	
	\vspace{0.4em}
	\textbf{Correct Answer: B}
	
	\textbf{Explanation:}
	A, C, D are incorrect. Incomplete approaches.
	B is correct. Comprehensive validation includes all four elements.
	
	\textbf{Learning Objective:} Discuss model validation and its importance.
	\textbf{Reference:} Module 5
	
	% =====================================================
	\subsection*{Question 98}
	Which is the best response to a model validation finding of material
	concern?
	
	A. Ignore the finding if the model is already in production.\\
	B. Escalate, remediate, and document before continued use.\\
	C. Deploy the model anyway with a plan to fix issues later.\\
	D. Ask the validation team to soften the finding.
	
	\vspace{0.4em}
	\textbf{Correct Answer: B}
	
	\textbf{Explanation:}
	A, C, D are incorrect. Unethical and non-compliant.
	B is correct. Escalate, remediate, document.
	
	\textbf{Learning Objective:} Discuss model validation and its importance.
	\textbf{Reference:} Module 5
	
	% =====================================================
	\subsection*{Question 99}
	Which best describes model drift?
	
	A. The model's underlying source code is updated to a new version.\\
	B. The statistical relationship between inputs and outputs changes over
	time.\\
	C. The hardware infrastructure hosting the model is upgraded.\\
	D. The vendor supporting the model is replaced by another vendor.
	
	\vspace{0.4em}
	\textbf{Correct Answer: B}
	
	\textbf{Explanation:}
	A, C, D are incorrect. Not the definition of drift.
	B is correct. Drift = change in input-output relationship.
	
	\textbf{Learning Objective:} Describe AI/ML model review framework.
	\textbf{Reference:} Module 5
	
	% =====================================================
	\subsection*{Question 100}
	Which best describes champion-challenger testing?
	
	A. Comparing a current production model with an alternative before
	switching.\\
	B. Comparing two identical models to verify they produce the same output.\\
	C. Testing a model without any data to verify robustness.\\
	D. Ignoring model performance in favor of qualitative assessments.
	
	\vspace{0.4em}
	\textbf{Correct Answer: A}
	
	\textbf{Explanation:}
	A is correct. Champion-challenger compares current vs.\ alternative.
	B, C, D are incorrect. Not the definition.
	
	\textbf{Learning Objective:} Describe steps in model adaptation.
	\textbf{Reference:} Module 5
	
	\newpage
	\section*{Master Answer Key}
	
	\begin{center}
		\begin{tabular}{|c|c|c|c|c|c|c|c|c|c|}
			\hline
			Q & A & Q & A & Q & A & Q & A & Q & A \\
			\hline
			1 & A & 21 & B & 41 & B & 61 & A & 81 & D \\
			2 & B & 22 & A & 42 & B & 62 & A & 82 & B \\
			3 & C & 23 & B & 43 & B & 63 & A & 83 & B \\
			4 & B & 24 & A & 44 & C & 64 & A & 84 & B \\
			5 & D & 25 & B & 45 & B & 65 & A & 85 & B \\
			6 & A & 26 & B & 46 & A & 66 & B & 86 & B \\
			7 & C & 27 & B & 47 & B & 67 & A & 87 & A \\
			8 & B & 28 & B & 48 & A & 68 & B & 88 & B \\
			9 & B & 29 & B & 49 & B & 69 & B & 89 & B \\
			10 & B & 30 & B & 50 & B & 70 & A & 90 & B \\
			11 & B & 31 & B & 51 & A & 71 & B & 91 & B \\
			12 & B & 32 & B & 52 & B & 72 & D & 92 & A \\
			13 & C & 33 & B & 53 & A & 73 & D & 93 & B \\
			14 & B & 34 & B & 54 & A & 74 & B & 94 & B \\
			15 & B & 35 & B & 55 & B & 75 & B & 95 & B \\
			16 & A & 36 & A & 56 & B & 76 & D & 96 & A \\
			17 & B & 37 & B & 57 & B & 77 & D & 97 & B \\
			18 & A & 38 & A & 58 & B & 78 & B & 98 & B \\
			19 & B & 39 & C & 59 & B & 79 & B & 99 & B \\
			20 & C & 40 & D & 60 & A & 80 & B & 100 & A \\
			\hline
		\end{tabular}
	\end{center}
	
\end{document}